\documentclass[twocolumn]{article}

\usepackage[margin=1.3cm]{geometry}
\usepackage[utf8]{inputenc}
\usepackage[T2A]{fontenc}
\usepackage[russian]{babel} 
\usepackage{amsmath, amssymb, amsfonts} 
\usepackage{graphicx} 
\usepackage{booktabs} 
\usepackage{algorithm} 
\usepackage{algorithmic} 
\usepackage[hidelinks]{hyperref}
\usepackage{xcolor}
\usepackage{csquotes}
\usepackage[backend=biber, style=numeric]{biblatex}
\addbibresource{cites.bib}

\title{\textbf{Multi-branch CNN based with multi-head attention multi-class brain disease classification using MPCA reduced dimensionality MRI scans over orthogonal planes}}
\author{Берлет Максим Валерьевич \texttt{berletmaximq@gmail.com} \\ Филатов Антон Юрьевич \texttt{ant.filatov@gmail.com}}
\date{}

\begin{document}

\maketitle

\begin{abstract}

Своевременная диагностика заболеваний, которые определяются по данным магнитно-резонансной томографии (МРТ), имеет решающее значение для выбора оптимальной тактики лечения. 
Автоматизация данного процесса обусловлена необходимостью снижения нагрузки на специалистов. 
Целью данной работы является разработка алгоритма классификации пациентов по МРТ-изображениям, эффективного в условиях ограниченных вычислительных ресурсов и малого объёма размеченных данных.
Предложен подход, основанный на декомпозиции трёхмерного МРТ-тензора по ортогональным плоскостям(Multilinear Principal Component Analysis), извлечении признаков с помощью многоветвистой сверточной нейросети и последующем позднем слиянии (late fusion) признаков с применением механизма многоголового внимания (multi-head attention) 
для адаптивного взвешивания информационных потоков. Оценка проводилась на выборке из 247 снимков пациентов методом 5-кратной кросс-валидации. В качестве метрик использованы точность (Accuracy) и F1-мера.
На валидационных фолдах модель продемонстрировала Accuracy = 1,00 и F1 = 1,00. Высокие показатели согласуются с характеристиками используемой выборки и подтверждают внутреннюю согласованность алгоритма.

Ключевые слова --- Multilinear Principal Component Analysis; Многоветвистая сверточная нейросеть; Многоголовое внимание; Позднее слияние признаков; Классификация; МРТ;

\end{abstract}

\section{Введение}
Нейродегенеративные заболевания представляют собой группу прогрессирующих патологий центральной нервной системы, характеризующихся постепенной гибелью нейронов и нарушением функций головного мозга. 
К таким заболеваниям относятся болезнь Альцгеймера, болезнь Паркинсона, рассеянный склероз и другие. 
В связи с глобальным старением населения их распространённость неуклонно растёт, что делает раннюю диагностику одной из наиболее актуальных задач современной медицины.

Магнитно-резонансная томография (МРТ) является золотым стандартом неинвазивной диагностики данных патологий благодаря высокому тканевому контрасту и возможности получения трёхмерных изображений. 
Однако интерпретация МРТ-исследований требует высокой квалификации и значительных временных затрат, что приводит к дефициту кадров и задержкам в постановке диагноза.

В последние годы активно развиваются методы автоматического анализа медицинских изображений на основе глубокого обучения. 
Большинство современных подходов демонстрируют высокую точность при использовании больших размеченных датасетов и мощных графических ускорителей. 
Тем не менее, в реальной клинической практике часто встречаются ограничения по объёму размеченных данных и вычислительным ресурсам.

Существующие методы можно разделить на две основные группы. К первой относятся подходы, работающие напрямую с трёхмерными данными (3D-CNN, Transformer-based модели). Они обычно показывают наилучшее качество, но требуют значительных вычислительных ресурсов. 
Ко второй группе относятся методы, основанные на обработке отдельных срезов (2D-CNN). Такие подходы менее ресурсоемкие, однако часто теряют важную пространственную информацию.

В связи с этим остаётся актуальной задача разработки моделей классификации МРТ-изображений, 
которые сочетали бы высокую точность с приемлемыми требованиями к вычислительным ресурсам и малому объёму обучающих данных.

Целью настоящей работы является разработка эффективного алгоритма классификации пациентов по МРТ-изображениям головного мозга в условиях ограниченных вычислительных ресурсов и 
малого размера размеченной выборки. 

Для достижения цели предлагается комбинированный подход, включающий:
\begin{itemize}
    \item декомпозицию трёхмерного МРТ-тензора по ортогональным плоскостям с использованием Multilinear Principal Component Analysis (MPCA);
    \item извлечение признаков с помощью многоветвистой свёрточной нейронной сети;
    \item позднее слияние признаков (late fusion) с применением механизма многоголового внимания (multi-head attention) для адаптивного взвешивания вкладов различных проекций.
\end{itemize}

Основной вклад работы заключается в разработке ресурсоэффективной архитектуры, демонстрирующей высокую обобщающую способность при крайне ограниченном объёме размеченных данных.

\section{Аналитический обзор}
Большинство современных работ в области классификации МРТ-изображений можно разделить на два направления: работы, которые анализируют срезы МРТ снимков в одной плоскости, а также 
работы, которые напрямую работают со всем МРТ тензором сразу. Основным недостатком первых работ является отсутствие учета пространственной информации с других срезов, 
а вторые работы, как правило, требуют значительных вычислительных ресурсов. 

В работе \cite{1} представлена легковесная двухпутевая сверточная нейросеть для мультиклассовой классификации опухолей головного мозга по предобработанным аксиальным срезам МРТ. 
Архитектура сети включает две параллельные ветви с различными размерами сверточных ядер: меньшее ядро ориентировано на извлечение локальных текстурных признаков, в 
то время как крупное ядро улавливает глобальный пространственный контекст. Извлеченные карты признаков конкатенируются и преобразуются в одномерный вектор. 
Ключевой особенностью подхода является интеграция модуля трехэтапного статистического отбора, который изолирует $K$ наиболее информативных признаков, 
существенно снижая размерность входного пространства последующего полносвязного слоя. К достоинствам метода относится высокая обобщающая способность (0.96 Accuracy) 
при минимальных вычислительных затратах (всего 861 545 параметров). 
Однако существенным ограничением работы является использование МРТ-данных только в одной проекции, что потенциально снижает устойчивость модели к пространственным вариациям опухолей в трехмерном объеме.

Метод основан на выделении признаков из области интереса (ROI), включающих геометрические, статистические и пространственно-текстурные характеристики. 
С целью снижения размерности пространства признаков авторы применили метод отбора на основе корреляционного анализа (CFS), что позволило сократить вектор признаков 
до 11 наиболее информативных компонентов. Сформированные векторы подавались на вход различных классификаторов, включая классический многослойный перцептрон (MLP), 
рекуррентную нейронную сеть (RNN) и ансамблевые методы машинного обучения. Ключевым достоинством данной работы является успешное решение задачи многоклассовой классификации
(дифференциация 6 типов новообразований, тогда как большинство существующих исследований сфокусированы на бинарной задаче), а также достижение высокой точности 
(Accuracy до 0.986) при использовании алгоритма Random Committee. Однако существенным недостатком предложенного конвейера является его строгая зависимость от качества 
предварительной бинарной сегментации (алгоритма ER-BHS): любая ошибка на этапе выделения контуров ROI приводит к искажению вычисляемых признаков. 
Кроме того, предложенный подход ограничен анализом лишь одиночных двухмерных (2D) срезов, что не позволяет учитывать полную трехмерную пространственную структуру опухоли.

Использование трансформеров для анализа медицинских изображений рассматривается в работе \cite{3}, где авторы модифицировали архитектуру ViT 
(модель FTVT-l16), адаптировав её под мультиклассовую классификацию МРТ-сканов с accuracy до 0.987. 
Однако, несмотря на высокие метрики, предложенное решение обладает рядом ограничений. 
Во-первых, анализ ограничен лишь одной плоскостью сканирования. 
Во-вторых, архитектура ViT остается требовательной к вычислительной инфраструктуре: для её обучения авторам потребовалось привлечение аппаратных ресурсов платформы Kaggle (две GPU NVIDIA T4).

В работе \cite{4} авторы предложили многоветвистую архитектуру на основе блоков Inception для мультиклассовой классификации МРТ-изображений головного мозга. 
Благодаря параллельным ветвям с различными размерами ядер свертки ($1\times1$, $3\times3$, $5\times5$), модель эффективно адаптируется к мультимасштабным признакам 
новообразований, демонстрируя высокую точность (Accuracy до 0.993) при пятикратной кросс-валидации. Тем не менее, подход имеет ряд ограничений: 
классификация проводится на уровне изолированных 2D-срезов, что приводит к потере пространственной 3D-информации МРТ-исследования. 
Кроме того, обучение сложной многоветвистой архитектуры с нуля на относительно небольшом объеме данных (2870 изображений) повышает риск переобучения и 
ограничивает обобщающую способность модели.

В работе \cite{5} авторы предложили классическую Feed-Forward сверточную нейронную сеть (MCCNN) для решения задачи многоклассовой классификации опухолей головного мозга. 
Предложенная архитектура продемонстрировала высокую точность (Accuracy) как на предобработанных данных после удаления черепа (до 0.99), так и на «сырых» снимках, 
содержащих костные структуры (до 0.96). Существенным преимуществом данного подхода является высокая параметрическая эффективность модели (всего 13 млн параметров), 
что снижает требования к вычислительным ресурсам. Однако ключевым ограничением работы остается ее двумерный характер: модель классифицирует отдельные срезы изолированно, 
игнорируя трехмерный пространственный контекст реального МРТ-сканирования.

Несмотря на значительный прогресс, большинство существующих решений либо требуют больших вычислительных ресурсов и объёмов данных, 
либо ограничиваются анализом одной плоскости сканирования. В отличие от них, предлагаемый в данной работе подход не требует значительных вычислительных ресурсов и больших объемов данных засчет использования предобученных моделей, 
притом максимально эффективно анализирует сразу все 3 плоскости МРТ одного снимка и дает ответ основанный на всей информации.

\section{Метод}

\subsection{Предобработка данных}
Описание предобработки данных.

\subsection{Архитектура модели}
Описание архитектуры.

\subsection{Алгоритм обучения}
Описание обучения.

\section{Эксперименты и результаты}
Описание экспериментов и таблиц.

\subsection{Описание датасета}

\subsection{Протокол эксперимента}

\section{Заключение}
Выводы по работе.

\printbibliography

\end{document}